% Section 7: Texas in context. Balanced: the state's program is small but fully
% spent, which is more than several peers can say.
\section{Texas commits little to music, and awarded all of it in the one round it has published}
\label{sec:state}

Austin's fund is not the only public money aimed at this workforce, and the
state's works differently. The Texas Music, Film, Television and Multimedia
Office, an office in the Governor's office of Economic Development and Tourism,
administers both the music rebate and the state's moving-image incentive under
Chapter 485 of the Government Code.\footnote{Government Code Chapter 485
places the Music, Film, Television and Multimedia Office inside the Governor's
Economic Development and Tourism division and gives it authority to administer
both incentive programs. That statutory location is why the same office serves
as the state actor across the film-and-television, film-workforce, and music
programs compared below.}
Placing those dedications beside one another shows what the state has committed
to each. This section then compares Texas with four peer states (Georgia, New
York, Tennessee, Louisiana) on music money awarded, claimed or appropriated
rather than merely authorised, sets the Texas arts economy beside the national
aggregate, and ends with the federal pandemic relief program that was the
largest single intervention of the period.

\subsection{The same state office commits far more to film than to music}

\begin{figure}[htbp]
\centering
\includegraphics[width=0.86\textwidth]{fig15_film_vs_music.png}
\caption{Texas dedications for the 2026--27 biennium, nominal dollars. Sources:
General Appropriations Act riders; Senate Bill 22, 89th Legislature, 2025. How
to read it: all three amounts are dedicated for fiscal 2026--27, so they are
already in comparable nominal terms and no inflation adjustment is needed.
Caveat: the film workforce line
is a one-time fiscal 2026 grant to universities rather than an ongoing
per-biennium dedication, it has no music counterpart, and it is never summed
with the other two.}
\label{fig:filmmusic}
\end{figure}

The 89th Texas Legislature dedicated \NTxFilmIncentiveBiennium\ to the
moving-image incentive for the 2026--27 biennium through Senate Bill 22, and
dedicated \NTxMusicTmirBiennium\ to the Texas Music Incubator Rebate over the
same two years under the earlier statute that created it. The ratio between the
two dedications, both nominal dollars for the same fiscal biennium, is
\NTxFilmMusicRatio.\footnote{The moving-image incentive appears at Texas Tax
Code section 151.801(g), added by Senate Bill 22 of the 89th Legislature
(2025), which requires the Comptroller to deposit \$300{,}000{,}000 per
biennium into the fund and sunsets on 1 September 2035. The music rebate
appears at Tax Code sections 183.023(c) and 151.801(f), added by Senate Bill
609 of the 87th Legislature (2021), which together dedicate
\$10{,}100{,}000 a year to the rebate account beginning with the 2024--25
biennium; Rider 39 or 40 of the General Appropriations Act appropriates that
amount to the Music Office each biennium. Because both figures are nominal
dollars for fiscal 2026--27, no inflation adjustment applies to their ratio.
Statutory materials at
\url{https://capitol.texas.gov/BillLookup/History.aspx?LegSess=89R\&Bill=SB22}
and the Comptroller's Manual of Accounts entry for the music account at
\url{https://fmcpa.cpa.state.tx.us/fiscalmoa/fund.jsp?num=5193}.}
The Texas Music Incubator Rebate, unlike its peer credits in other states,
returns a portion of mixed-beverage and sales taxes that a qualifying venue or
promoter has already remitted to the Comptroller, so the subsidy scales with a
venue's alcohol sales rather than with what the venue pays performers. The
authorising statute requires that a written artist contract exist but sets no
minimum compensation.\footnote{The Legislature created the rebate in Senate
Bill 609 in 2021 and left it unfunded until 2023, and the first application
window opened on 1 September 2023.} The same 2025 legislative session that
increased film also authorised a \NTxFilmWorkforcePilotOnetime\ Higher
Education Film Workforce Pilot Program under Rider 45 of the 2026--27 General
Appropriations Act, a one-time fiscal 2026 grant to universities for film and
media workforce training with no music counterpart.\footnote{Rider 45 restricts
the \$30{,}000{,}000 of General Revenue to eligible general academic teaching
institutions and to film and media workforce activities specifically, and
prohibits its use for any other purpose. Registered as one-time so it is never
summed with the recurring biennial moving-image or music dedications.}

\subsection{A statutory cap is not a measure of support}

\begin{figure}[htbp]
\centering
\includegraphics[width=\textwidth]{fig16_state_percapita.png}
\caption{State music incentives per resident, authorised against actually
awarded, claimed or appropriated, five states, 2025 dollars, with general state arts agency funding
shown separately below. Sources: state statutes and agency reports; National
Assembly of State Arts Agencies; Census Bureau population estimates for July
2025. How to read it: the upper panel shows what each state's music-specific
program allows against what it paid, and the lower panel is general arts
funding, a different measure that does not add to the upper one. Caveats:
the New York authorised figure is a proxy derived from a budget increase rather
than a flat annual cap; the Tennessee bars combine film, television and music
scoring with no music-only breakout available, so they are an upper bound.}
\label{fig:statepercap}
\end{figure}

The Georgia General Assembly authorised up to \$15 million a year for a music
tax credit under O.C.G.A.\ section 48-7-40.33, and the Georgia Department of
Revenue paid out nothing over the credit's five-year active life. The Georgia
Department of Audits and Accounts, in a December 2023 evaluation, found that
the department received six applications and approved none, because each
lacked the required pre-certification from the Georgia Department of Economic
Development.\footnote{Georgia Department of Audits and Accounts and Georgia
Southern University CBAER, \textit{Georgia Music Tax Credit: Economic and
Fiscal Impact Analysis}, report dated 8 December 2023, published 14 December
2023, at
\url{https://www.audits.ga.gov/ReportSearch/download/30448}. The state's
published Tax Expenditure Reports listed the credit as ``less than
\$1 million,'' a notation that concealed a true value of zero.} On the
statutory ceiling alone, Georgia's music support looks larger than Texas's;
on dollars actually paid, Georgia is at
\NMusicDisbursedPercapitaGeorgia\ per resident.

On money awarded, claimed or appropriated per resident rather than authorised,
Texas is at \NMusicDisbursedPercapitaTexas\ against New York's
\NMusicDisbursedPercapitaNewyork\ and Louisiana's
\NMusicDisbursedPercapitaLouisiana. Tennessee's
\NMusicDisbursedPercapitaTennessee\ combines film, television and music
scoring together in the Tennessee state budget's single Film and Television
Incentive Fund line, and so cannot be ranked against a music-only figure.
Georgia paid nothing at all. Among the states here whose figures are
music-specific, Texas's dedication is the smallest per resident of those that
paid anything.\footnote{The per-resident figures divide two different
dollar quantities by the same population base. Authorised per resident is
\(a_s = A_s\,d_{y(A_s)} / N_s\) and disbursed per resident is
\(b_s = B_s\,d_{y(B_s)} / N_s\), where \(A_s\) is the annual amount state
\(s\) may award under the statute or dedication behind its music-specific
program, \(B_s\) is the amount the state actually paid or certified,
\(d_y = C_{2025}/C_y\) carries each figure from its own nominal year to 2025
dollars using the national CPI-U, and \(N_s\) is the state's resident
population on 1 July 2025 from the Census Bureau's Vintage 2025 state
population totals at
\url{https://www2.census.gov/programs-surveys/popest/datasets/2020-2025/state/totals/NST-EST2025-ALLDATA.csv}.
The two figures can differ by the entire amount, because an authorisation is a
ceiling written into law and a disbursement is money that left the treasury.
Georgia shows the gap at its widest, with a positive ceiling and nothing paid.
Where a state's authorisation is not a flat annual cap, the figure is a proxy:
the New York number is the most recent one-year budget increase to a
multi-year lifetime cap under Tax Law section 24-c, and the Tennessee figure
combines film, television and music scoring because the Tennessee Department
of Finance and Administration reports no music-only breakout at line 330.17 of
the state budget document, so both are upper bounds. Three cautions apply to
the second column. Only the Texas and Georgia figures are money awarded that
this study could verify against agency award files; the New York, Tennessee
and Louisiana figures are respectively an annualised award total, a
combined-fund appropriation and a claimed credit, so the column is not a
like-for-like disbursement measure. The New York figure annualises
\$308 million in credits awarded between August 2022 and February 2026 (three
and a half years) into a one-year equivalent, against a single Texas rebate
round. The Louisiana figure combines the Sound Recording Investor Tax Credit
at Louisiana Revised Statute 47:6023 with the Musical and Theatrical
Production Credit at R.S. 47:6034; the Sound Recording credit alone is about
\$0.01 per resident, well below Texas. These are accounting ratios built from
published totals, not estimates, and no sampling error attaches to them.
Program-page URLs and the underlying budget documents are consolidated in the
technical appendix under each state name.}

The comparison misses how differently the Texas program is built. In the one
rebate round with a published total, fiscal 2026, the Music Office awarded its
full annual dedication, which no peer state here did. The fiscal 2025 round
total has not been published, and the first round's total is internally
inconsistent, so ``full spending'' describes one documented round rather than a
program-level record. Texas is also the only state compared here that
subsidises operating music venues at all: the peer credits fund production
companies mounting touring or Broadway-style shows, a different instrument
aimed at a different beneficiary.\footnote{The five-state comparison is
therefore not a league table. New York's musical and theatrical production
credit at Tax Law section 24-c subsidises productions; Tennessee's
music-related incentive is a film-scoring subsidy under the 2018 Visual
Content Modernization Act, and the Tennessee Entertainment Commission's
program guidelines expressly declare music videos ineligible. Louisiana's two
credits (sound recording investment and musical and theatrical production)
subsidise recording sessions and touring productions respectively. Only the
Texas rebate pays the operating room itself.}

\subsection{Texas has a small arts economy relative to its size}

\begin{figure}[htbp]
\centering
\includegraphics[width=\textwidth]{fig17_bea_arts.png}
\caption{Arts and cultural production as a share of state economic output, eight
states, 2023; and Texas value added for performing arts companies against
independent artists, writers and performers, 2001--2023. Source: Bureau of
Economic Analysis Arts and Cultural Production Satellite Account, nominal
dollars. How to read it: the upper panel compares the eight benchmark states
only, not all fifty. Caveat: this account has no metro-level detail, so it says
nothing about Austin specifically.}
\label{fig:beaarts}
\end{figure}

The Bureau of Economic Analysis publishes the Arts and Cultural Production
Satellite Account, an annual accounting series that measures the value added
by a defined set of arts and cultural industries against total gross domestic
product; the account is the only official U.S.\ statistic that isolates the
economic size of the arts sector at the state
level.\footnote{The BEA account, generally abbreviated ACPSA, uses production
categories built from the North American Industry Classification System and
covers roughly 35 industries from performing arts and museums to
architecture, publishing and broadcasting. Its state file at
\url{https://www.bea.gov/data/special-topics/arts-and-culture} is a
top-down allocation of national commodity totals to states and does not
resolve below the state level, so it cannot describe Austin specifically. The
2023 estimates released on 2 April 2025 in BEA news release 25-13 are the
latest available. The share it produces is the appropriate measure of an
industry's size within a state economy because it deflates the industry to
that state's own economic base, which lets a comparison across states control
for the size of the state.} Arts and cultural production is
\NBeaGdpShareTxTwoZeroTwoThree\ of Texas economic output in 2023, against
\NBeaGdpShareUsAvgTwoZeroTwoThree\ for the United States as a whole. The
United States figure is the national aggregate ratio (all U.S.\ arts value
added divided by total U.S.\ GDP) rather than the mean of the fifty state
shares.\footnote{The share for state \(s\) is
\(100 \times V_s / Y_s\), where \(V_s\) is nominal value added by arts and
cultural production industries in state \(s\) and \(Y_s\) is that state's
nominal total gross domestic product, both in 2023. A ratio of two same-year
dollar figures needs no inflation adjustment, because deflating numerator and
denominator by the same factor leaves the ratio unchanged. Among the eight
benchmark states plotted (California, Florida, Georgia, Louisiana, New York,
Tennessee, Texas and Washington), only Louisiana ranks lower than Texas, at
\NBeaGdpShareLaTwoZeroTwoThree. Nationally, Texas ranks
\NBeaGdpShareTxNationalRank\ across all 50 states and the District of
Columbia; this ranking is taken from BEA's own release rather than re-derived
here, since this study's extract holds only the eight benchmark states and
does not permit an independent national ranking.} The composition inside the
Texas total has also shifted. Value added by independent artists, writers and
performers in Texas first overtook that of performing arts companies in
\NBeaFirstCrossoverYearTx, fell back behind performing arts companies for two
years, then passed it for good in \NBeaCrossoverYearTx\ and has led every year
since. The same movement from institutions toward individuals that
Section~\ref{sec:payroll} found in payroll employment appears here in output.

The Bureau of Economic Analysis posted a notice in February 2026 stating that
it will no longer regularly produce the Arts and Cultural Production
Satellite Account, and so the April 2025 release covering 2023 is the last
year available and no further annual update is
scheduled.\footnote{BEA site text, ``BEA will no longer regularly produce these
statistics,'' February 2026, retrieved 1 August 2026 from
\url{https://www.bea.gov/data/special-topics/arts-and-culture}, cross-checked
against the April 2025 release (BEA news release 25-13). BEA has not ruled
out a future non-regular release, so this is a discontinued schedule rather
than a closed series.}

\subsection{Federal relief was the largest intervention, and it was one-off}

\begin{figure}[htbp]
\centering
\includegraphics[width=0.86\textwidth]{fig18_svog_metro.png}
\caption{Shuttered Venue Operators Grant awards by metro area. Source: Small
Business Administration award file, July 2022 vintage, the last published. How
to read it: bars are total awarded dollars, nominal. Caveats: metro groupings are built
from recipient city names rather than official metropolitan boundaries, so they
are approximate and will differ from other researchers' groupings for the same
metros; recipient city is the filer's mailing address, so a multi-site chain
books grants for out-of-state locations to its headquarters city; the Nashville
figure covers the consolidated city only while the two Texas figures use
multi-place lists; and the Austin total covers every eligible entity type, of
which motion picture theatre operators account for \$109.8 million across 24
recipients against \$104.9 million to 80 live venue operators and promoters.}
\label{fig:svog}
\end{figure}

Congress created the Shuttered Venue Operators Grant program in December 2020
to compensate live venue operators, promoters, theatrical producers and
motion picture theatre operators for revenue lost during the pandemic
closures; the Small Business Administration administered the program and
published the award file used here.\footnote{The Small Business
Administration award file used here, ``Awards as of 07-05-2022,'' is the
last public release SBA has posted to its open data portal at
\url{https://data.sba.gov/sites/default/files/distribution/SBA-ODA-2022-09-001/awards-as-of-7-5-22.xlsx};
it holds 13{,}011 recipient records nationally. The program's authorising
statute, section 324 of the Economic Aid to Hard-Hit Small Businesses,
Nonprofits, and Venues Act (division N of the Consolidated Appropriations
Act, 2021), capped awards at ten million dollars per entity and gave
priority to operators reporting the deepest revenue drops. Documentation at
\url{https://www.congress.gov/crs_external_products/R/PDF/R46689/R46689.5.pdf}
and \url{https://www.sba.gov/article/2025/01/15/new-report-shows-consequential-impacts-sba-pandemic-relief}.}
Austin-area recipients received \NSvogAustinTotal\ in federal shuttered-venue
grants, against \NSvogHoustonAreaTotal\ for this study's Houston-area
grouping, \NSvogNashvilleTotal\ for Nashville and \NSvogDfwAreaTotal\ for
Dallas-Fort Worth. The Austin and Houston figures should not be ranked
against one another. Three awards inside the Austin group, together worth
\$16.5 million, are filed at Round Rock and Austin corporate addresses but
name cinemas in Iowa, Indiana and Virginia; removing them alone puts Austin
below Houston. Money was concentrated: the top ten Austin recipients
received \NSvogAustinTopOneZeroShare\ of Austin dollars, and
\NSvogAustinCapHits\ reached the program's \$10 million per-entity
cap.\footnote{The concentration measure is
\(100 \times \sum_{j=1}^{10} G_{(j)} \big/ \sum_{j} G_j\), where \(G_j\) is the
awarded amount for recipient \(j\) in the Austin group and \(G_{(1)} \ge
G_{(2)} \ge \dots\) are the same awards sorted from largest to smallest, so
the numerator is the ten largest awards and the denominator is every Austin
award. Award dollars are nominal and are not deflated, because all awards
fall in the same program period. The Austin group is built from recipient
city names rather than official metropolitan boundaries, so the denominator
is this study's grouping and will differ from another researcher's using a
different city list.}

The event study around those awards illustrates why relief programs resist
causal evaluation. In the six months before their award, recipient venues
were running about 39 percent below comparison venues on log monthly
mixed-beverage receipts, with a 95 percent confidence interval that stretched
from about 58 percent below to 12 percent below. The half-year before that is
estimated too imprecisely to distinguish from no gap. Recipients then
converged to no measurable difference from comparison venues after the award.
The pre-award trough looks like the shape of a program effect, but SVOG
eligibility required documenting a revenue drop, so the pre-award trough is
the selection rule made visible in the data rather than an independent
finding of harm.\footnote{The event study runs on half-year blocks measured
from each venue's own award month:
\[
\ln R^{\mathrm{real}}_{it}
  = \alpha_i + \tau_t
  + \sum_{b \neq b_0} \beta_b\,
    \bigl(\mathbf{1}\{\mathrm{block}_{it} = b\} \times \mathrm{SVOG}_i\bigr)
  + \varepsilon_{it},
\]
where \(R^{\mathrm{real}}_{it}\) is venue \(i\)'s monthly receipts in 2025
dollars from the Texas Comptroller Mixed Beverage Gross Receipts dataset
(naix-2893 on the Texas Open Data Portal at
\url{https://data.texas.gov/dataset/Mixed-Beverage-Gross-Receipts/naix-2893}),
\(\mathrm{SVOG}_i\) is one for a matched relief recipient,
\(\mathrm{block}_{it} = \lfloor (t - a_i)/6 \rfloor\) places month \(t\) in a
half-year block relative to that venue's own award month \(a_i\), clipped at
nine blocks before and eight after so the end blocks hold several months for
every award cohort rather than a sliver, and \(\alpha_i\) and \(\tau_t\) are
venue and calendar-month fixed effects respectively. The omitted block
\(b_0\) covers 30 to 25 months before the award, which for the typical award
month is the first half of 2019; the conventional base of the period just
before treatment would fall in the deepest part of the pandemic and would
leave every other coefficient unreadable. Venues that never received an
award are assigned the base block, so they contribute to the fixed effects
and the common time path rather than to any \(\beta_b\). Standard errors are
clustered by venue, and the largest ticketed rooms are excluded. The
p-value quoted below is a joint Wald test that the block coefficients
closing before January 2020 are all zero (the blocks that predate the
pandemic shock for every award month in the data). The pre-pandemic trends
for the two groups are not distinguishable
(\(p = \)~\NQuasiSvogPretrendPPrepandemic), the strongest parallel-trends
evidence in this study, and even so the design still cannot support a causal
claim about the awards themselves. This study matched
\NQuasiSvogMatchedVenues\ Austin awards to the curated venue panel and
dropped three matches as uncertain, with the reasons recorded in the
analysis files. The awards all arrived within a three-month window in
2021, so relative time and calendar time are close to relabellings of one
another.}

The largest single intervention in this ecosystem over the period studied
was federal, one-off, and directed at venue and cinema operators rather than
at musicians. Of the Austin-area total, \$104.9 million went to live venue
operators and promoters (80 recipients) and \$109.8 million to motion
picture theatre operators (24 recipients), with the remainder going to
performing arts organisations, talent representatives, museums and
theatrical producers.

Sections~\ref{sec:earnings} through \ref{sec:state} have set out what the
public record measures. Section~\ref{sec:lessons} pairs each finding with the
option it points toward and the precedent this study located, so a reader can
judge the inference for themselves before Section~\ref{sec:next} names the
additions that would close the measurement gaps this analysis ran into.
